\stfhIndex{I091}{Ziyouhuanlukongjian}{自由环路空间（free loop space）}
\section{\texorpdfstring{自由环路空间及其 $S^1$-作用}{自由环路空间及其 S¹-作用}}
\label{sec:free-loop-space-circle-action}

设 $M$ 是一般光滑流形，不必是辛流形。记
$\mathcal L(M):=\operatorname{Map}(S^1,M)$ 为\emph{自由环路空间}，即光滑映射
\[
  \gamma:S^1=\mathbb R/\mathbb Z\longrightarrow M
\]
的集合。我们强调，环路带有标记点 $0\in\mathbb R/\mathbb Z$；我们也常以单位区间
$[0,1]$ 参数化环路，并附加周期边界条件 $\gamma(0)=\gamma(1)$。
$\mathcal L(M)$ 有一个由可缩环路组成的特定连通分支，记作 $\mathcal L_0(M)$。
其\termUniversalCoveringSpace{}，记作 $\widetilde{\mathcal L}_0(M)$，可以写成丛
$\pi:\widetilde{\mathcal L}_0(M)\to\mathcal L_0(M)$，在
$\gamma\in\mathcal L_0(M)$ 上的纤维为
\[
 \{[\gamma,w]\mid w:D^2\to M,\ \partial w:=w|_{\partial D^2}=\gamma\},
\]
其中 $[\gamma,w]$ 表示满足 $\partial w=\gamma$ 的圆盘
$w:D^2\to M$ 相对于边界的同伦类。这里通过乘以 $2\pi$ 将 $\partial D^2$ 与
$S^1$ 认同。称这样的 $w$ 为 $\gamma$ 的\emph{填充圆盘}。

由纤维化 $\Omega_0(M)\to\mathcal L(M)\to M$ 可得
\[
 \pi_1(\mathcal L(M))\cong\pi_1(\Omega_0(M))\times\pi_1(M),
\]
其中 $\Omega_0(M)$ 是基环路空间，比如以常环路
$\widetilde y_0\in\mathcal L_0(M)$ 为基点，且 $y_0\in M$。回忆
$\pi_1(\Omega_0(M))=\pi_1(\Omega_0(M),\widetilde y_0)\cong\pi_2(M)$，便得到
\footnote{译者注：原书在这段纤维化中仍以 $\mathcal L_0(M)$ 表示基环路空间，
与前文同一符号所表示的可缩自由环路分支冲突；译文以常用记号 $\Omega_0(M)$ 区分
基环路空间，而式 (18.1.1) 仍照录原书的群同构。}
\footnote{译者注：在一般光滑流形这一假设下，评价纤维化通常只给出
$\pi_2(M)$ 被 $\pi_1(M)$ 的单值性作用扭曲的半直积，而不一定是式 (18.1.1) 所写的
直积；仅以填充圆盘为数据的覆盖也不一定是万有覆盖。原书未补充使这些断言成立的
额外假设，故正文照录并在此提示。}
\begin{equation}
 \pi_1(\mathcal L_0(M),\widetilde y_0)
 \cong\pi_2(M,y_0)\times\pi_1(M,y_0).
 \label{eq:18-free-loop-fundamental-group}
\end{equation}

与元素 $(A,a)\in\pi_2(M)\times\pi_1(M)$ 相伴的
$\widetilde{\mathcal L}_0(M)$ 的\termDeckTransformation{}可如下实现。以满足
$v(N)=y_0$ 的球面 $v:S^2\to M$ 表示 $A$，并取
$[\gamma,w]\in\widetilde{\mathcal L}_0(M)$。把区间 $[0,1]$ 看作
$D^2$ 中的子集 $[0,1]\times i\cdot\{0\}$，并以 $0$ 表示 $D^2$ 的原点。
任选道路 $\ell_1:[0,1]\to M$，使得
$\ell_1(0)=y_0$、$\ell_1(1)=w(0)$，把 $w(0)$ 与 $y_0$ 连接起来。
\footnote{译者注：源书印 $\ell_1(0)=w(0)$、$\ell_1(1)=0$，端点互换且基点误写为
$0$；译文补正为 $\ell_1(0)=y_0$、$\ell_1(1)=w(0)$，使拼接
$\ell_1\#w|_{[0,1]}\#\overline{\ell}_2$ 成为以 $y_0$ 为基点的环路。}
再选道路
$\ell_2:[0,1]\to M$，满足
\[
 \ell_2(0)=y_0,\qquad \ell_2(1)=\gamma(1),
\]
并考虑连接
\[
 \ell_1\#w|_{[0,1]}\#\overline{\ell}_2,
\]
它给出一个以 $y_0$ 为基点的环路。要求这个环路的同伦类就是给定的
$a\in\pi_1(M)$。特别地，$A\in\pi_2(M)$ 在 $[\gamma,w]$ 上的作用由“粘上一个
球面”的操作实现：
\begin{equation}
 (\gamma,w)\longmapsto(\gamma,w\#v),
 \label{eq:18-glue-sphere-action}
\end{equation}
其中模去上述 $\pi_1(M)$ 的作用，而 $v:S^2\to M$ 是表示 $A$ 的某个（因而任意一个）
球面。

$\mathcal L(M)$ 上有由定义域平移诱导的自然圆周作用
\begin{equation}
 \gamma\longmapsto\gamma\circ R_\varphi=\gamma(\,\cdot+\varphi),
 \label{eq:18-loop-rotation}
\end{equation}
其中 $R_\varphi:S^1\to S^1$ 定义为
\[
 R_\varphi(t)=t+\varphi,\qquad\varphi\in S^1.
\]
这一作用的无穷小生成元是 $\mathcal L(M)$ 上的向量场 $\mathbb X$：
\begin{equation}
 \mathbb X(\gamma)=\dot\gamma.
 \label{eq:18-circle-generator}
\end{equation}
该 $S^1$-作用的不动点集是常环路的集合
\[
 M\lhook\joinrel\longrightarrow\mathcal L_0(M).
\]
这一作用还提升为二元组 $(\gamma,w)$ 集合上的作用
\begin{equation}
 (\gamma,w)\longmapsto(\gamma\circ R_\varphi,w\circ R_\varphi),
 \label{eq:18-lifted-circle-action}
\end{equation}
其中再次用 $R_\varphi$ 表示复乘法所诱导的映射
\[
 R_\varphi:z\in D^2\subset\mathbb C\longmapsto e^{2\pi i\varphi}z.
\]

\begin{lemma}
\label{lem:18-fixed-points-lifted-circle}
$\widetilde{\mathcal L}_0(M)$ 上诱导的 $S^1$-作用，其不动点集由
$[\gamma_x,w_x]$ 组成，其中 $\gamma_x$ 和 $w_x$ 分别是像为 $x$ 的常环路和常圆盘。
\footnote{译者注：按前文的万有覆盖模型，常环路也可配以代表非零球面类的填充圆盘；
同样利用 $R_\varphi\simeq\mathrm{id}$，这类加帽常环路的同伦类似乎也为不动点。
原书本引理只列常圆盘，与第 18.2 节稍后
$\operatorname{Crit}\mathcal A_0=\bigsqcup_{g\in\Gamma}g\cdot M$ 的描述亦有张力；
因缺少进一步约定，译文照录本引理并在此注明。}
\end{lemma}

\begin{proof}
定义域旋转 $R_\varphi$ 在常环路 $\gamma\equiv x\in M$ 上作用平凡。另一方面，
它在二元组 $(\gamma_x,w_x)$ 的同伦类上也作用平凡，因为 $R_\varphi$ 同伦于恒等映射，
所以 $(\gamma_x,w_x)$ 与
$(\gamma_x\circ R_\varphi,w_x\circ R_\varphi)$ 彼此同伦。
\end{proof}
